\documentclass[11pt,a4paper]{article}
\usepackage[T1]{fontenc}
\usepackage[utf8]{inputenc}
\usepackage{amsmath,amssymb,amsthm}
\usepackage[margin=1in]{geometry}
\usepackage[colorlinks=true,linkcolor=blue,urlcolor=blue]{hyperref}
\theoremstyle{plain}
\newtheorem{theorem}{Theorem}
\newtheorem{lemma}[theorem]{Lemma}
\newtheorem{corollary}[theorem]{Corollary}
\theoremstyle{definition}
\newtheorem{remark}[theorem]{Remark}
\title{A proof of the conjectured recurrence for OEIS A245088}
\author{Adrian Perez Fontelles\\ \small Independent researcher}
\date{26 August 2026}
\begin{document}
\maketitle

\begin{abstract}
OEIS A245088 carries a conjectured linear recurrence with polynomial coefficients, of
order 5. The entry separately states a recurrence of order 3 that is not
labelled a conjecture. We prove that the first is a consequence of the second, so the
conjecture holds. The argument is a division in the Ore algebra
$\mathbb{Q}(n)[N]$, where $N$ is the shift $Nf(n)=f(n+1)$: writing $C$ and $P$ for the two
recurrences read as operators, we exhibit $Q$ with $C=QP$, whence
$C(a)=Q(P(a))=Q(0)=0$. No generating function is used. The division is exact
rational-function arithmetic, and the finitely many $n$ excluded by denominators of $Q$
are computed rather than ignored.
\end{abstract}

\noindent\small 2020 Mathematics Subject Classification. 11B37, 12H10, 33F10.\normalsize

\section{The sequence and the two recurrences}

OEIS A245088 is ``a(n) = (Sum\_\{k=0..n-1\} (8*k\textasciicircum{}2+12*k+5)*A244973(k))/n\textasciicircum{}2''. It has offset $1$ and begins
\[
a(1),\dots,a(8)\;=\;5, -5, -9, 115, -397, 85, 6625, -36181,\ \dots
\]
The entry carries the following comment, still recorded as a conjecture:
\begin{quote}\small
\textbf{Conjecture:} +n\textasciicircum{}2*(23656*n-2233)*(n-1)*a(n) +(n-1)*(23656*n\textasciicircum{}3+655645*n\textasciicircum{}2-2273556*n+1915600)*a(n-1) -2*(n-2) *(177048*n\textasciicircum{}3-3748011*n\textasciicircum{}2+13870329*n-15407140)*a(n-2) -2*(n-3) *(2020728*n\textasciicircum{}3-18138829*n\textasciicircum{}2+63885382*n-73802216)*a(n-3) -5*(n-4) *(782552*n\textasciicircum{}3-11993599*n\textasciicircum{}2+55520403*n-80051816)*a(n-4) +125*(66088*n-203507)*(n-5)*(n-4)\textasciicircum{}2*a(n-5)=0. - \emph{R. J. Mathar}, Jul 12 2014
\end{quote}
As of the ``Last modified'' line on the live entry (2026-05-30, revision 29) it is still
a conjecture, and no proof appears among the entry's links, formulas or programs.

The same entry also states, \emph{not} as a conjecture:
\begin{quote}\small
Recurrence (of order 3): (n-1)*n\textasciicircum{}2*(8*n\textasciicircum{}2 - 36*n + 41)*(8*n\textasciicircum{}2 - 20*n + 13)*a(n) =  -(n-1)*(8*n\textasciicircum{}2 - 36*n + 41)*(40*n\textasciicircum{}4 - 148*n\textasciicircum{}3 + 161*n\textasciicircum{}2 - 42*n + 4)*a(n-1) - (n-2)*(8*n\textasciicircum{}2 - 4*n + 1)*(152*n\textasciicircum{}4 - 1188*n\textasciicircum{}3 + 3427*n\textasciicircum{}2 - 4293*n + 1960)*a(n-2) + 25*(n-3)*(n-2)\textasciicircum{}2*(8*n\textasciicircum{}2 - 20*n + 13)*(8*n\textasciicircum{}2 - 4*n + 1)*a(n-3). - \emph{Vaclav Kotesovec}, Jul 13 2014
\end{quote}
This second relation is what the argument below rests on; it is taken as given, exactly as
a posted generating function would be. It was checked against the entry's own DATA before
being used, and holds on every published term it reaches.

\section{Recurrences as operators, and what division proves}

Let $N$ be the shift operator, $(Nf)(n)=f(n+1)$, acting on sequences. Polynomials in $n$
and $N$ generate the Ore algebra $\mathcal{A}=\mathbb{Q}(n)[N]$, in which multiplication is
associative and distributive but not commutative: moving $N$ past a coefficient shifts it,
\[
N\,q(n)\;=\;q(n+1)\,N .
\]
A recurrence $\sum_{i=0}^{r}p_{i}(n)\,a(n-i)=0$ becomes an element of $\mathcal{A}$ after
re-indexing $n\mapsto n+r$, namely
\[
L\;=\;\sum_{j=0}^{r}q_{j}(n)\,N^{j},\qquad q_{j}(n)=p_{r-j}(n+r),
\]
and the recurrence says exactly that $L$ annihilates the sequence, $L(a)=0$.

\begin{lemma}\label{lem:factor}
Let $P,C\in\mathcal{A}$ and suppose $C=QP$ for some $Q\in\mathcal{A}$. If $P(a)=0$ then
$C(a)=0$, at every $n$ where the coefficients of $Q$ are defined.
\end{lemma}

\begin{proof}
$C(a)=(QP)(a)=Q\bigl(P(a)\bigr)=Q(0)=0$. The only caveat is the one stated: a coefficient
of $Q$ is a rational function of $n$, so the step is valid at every $n$ that is not a pole
of one of them.
\end{proof}

\begin{lemma}\label{lem:div}
Right division works in $\mathcal{A}$: given $C$ and $P\neq0$, there are $Q,R\in\mathcal{A}$
with $C=QP+R$ and $\deg_{N}R<\deg_{N}P$.
\end{lemma}

\begin{proof}
Induct on $\deg_{N}C$. If $\deg_{N}C<\deg_{N}P$ take $Q=0$, $R=C$. Otherwise let
$d=\deg_{N}C-\deg_{N}P$ and let $c$, $p$ be the leading coefficients of $C$ and $P$. By the
commutation rule, the leading coefficient of $N^{d}P$ is $p(n+d)$, which is a nonzero
element of $\mathbb{Q}(n)$ and hence invertible there. Then
$C-\dfrac{c(n)}{p(n+d)}\,N^{d}P$ has strictly smaller degree in $N$, and induction applies.
\end{proof}

So the question ``does the conjecture follow from the stated recurrence?'' is decided by
running the division and looking at the remainder. Nothing is approximated: the
coefficients live in $\mathbb{Q}(n)$ and the arithmetic is exact.

\section{The computation for A245088}

Reading both relations as operators and dividing the conjectured one by the stated one
leaves remainder $0$, with quotient
\[
Q \;=\; \left(\frac{- 330440 n - 634665}{64 n^{4} + 576 n^{3} + 1920 n^{2} + 2808 n + 1525}\right) + \left(\frac{- 756992 n^{3} - 5263408 n^{2} - 11204200 n - 6603934}{512 n^{6} + 8448 n^{5} + 57152 n^{4} + 202752 n^{3} + 397640 n^{2} + 408804 n + 172325}\right)\,N + \left(\frac{23656 n + 116047}{64 n^{4} + 832 n^{3} + 4032 n^{2} + 8632 n + 6893}\right)\,N^{2} .
\]
Hence $C=QP$ in $\mathcal{A}$.

The coefficients of $Q$ are rational functions of $n$, and none of their poles is an integer, so Lemma~\ref{lem:factor} applies at every index.

\begin{theorem}
The conjectured recurrence
\[
n^{2} \left(n - 1\right) \left(23656 n - 2233\right)\,a(n) + \left(n - 1\right) \left(23656 n^{3} + 655645 n^{2} - 2273556 n + 1915600\right)\,a(n-1) -  2 \left(n - 2\right) \left(177048 n^{3} - 3748011 n^{2} + 13870329 n - 15407140\right)\,a(n-2) -  2 \left(n - 3\right) \left(2020728 n^{3} - 18138829 n^{2} + 63885382 n - 73802216\right)\,a(n-3) -  5 \left(n - 4\right) \left(782552 n^{3} - 11993599 n^{2} + 55520403 n - 80051816\right)\,a(n-4) + 125 \left(n - 5\right) \left(n - 4\right)^{2} \left(66088 n - 203507\right)\,a(n-5) \;=\;0
\]
holds for every $n\ge6$.
\end{theorem}

\begin{proof}
By Lemma~\ref{lem:factor} applied to the factorisation above, using that the stated
recurrence annihilates the sequence. No coefficient of $Q$ has an integer pole, so the lemma applies at every $n$.
\end{proof}

\section{Verification}

Every number above is an output of a script written separately from this note, and three
independent checks were run.

First, the factorisation was not assumed but computed, and then confirmed by multiplying
back: $QP$ was expanded in $\mathcal{A}$ and its coefficients cancelled against those of
$C$ to zero, as rational functions of $n$.

Second, the relation the argument rests on was checked against the entry rather than
trusted: the stated recurrence was evaluated on the published terms in exact integer
arithmetic and holds on every one it reaches.

Third, the conjecture itself was evaluated the same way, directly and without reference to
either operator: for each $n$ in range the sum $\sum_{i}p_{i}(n)a(n-i)$ was formed from the
entry's own values and found to vanish, for all 20 values of $n$ from
$n=6$ upward.

\begin{thebibliography}{9}
\bibitem{oeis} The OEIS Foundation, \emph{The On-Line Encyclopedia of Integer Sequences},
\url{https://oeis.org}, 2026. Sequence A245088.
\bibitem{ore} O.~Ore, \emph{Theory of non-commutative polynomials}, Ann. of Math. 34
(1933), 480--508.
\bibitem{kp} M.~Kauers and P.~Paule, \emph{The Concrete Tetrahedron}, Springer, 2011,
Chapter 7.
\bibitem{bronstein} M.~Bronstein and M.~Petkov\v{s}ek, \emph{An introduction to
pseudo-linear algebra}, Theoret. Comput. Sci. 157 (1996), 3--33.
\end{thebibliography}
\end{document}
